\clearpage
\VerdierPhysicalPageAnchor{0046}%
\VerdierSemanticAnchor{ch1-1-4-3-canonical}
この作用を\emph{標準作用}という。関手 $F$ 上に群 $\mathbb Z/2\mathbb Z$ の
標準作用を誘導する $H$-関手

\VerdierSemanticAnchor{ch1-1-4-3-eq3}
\begin{equation}
  (F,m)\colon(\mathcal C,H,T_f,c_f)
  \longrightarrow(\mathcal C',H,T'_f,c'_f)
  \tag{1.4.3.3}
\end{equation}
を調べる。

\VerdierSemanticAnchor{ch1-1-4-3-extension}
$\dot\alpha\in H^2(G,\mathbb Z/2\mathbb Z)$ を拡大

\VerdierSemanticAnchor{ch1-1-4-3-eq4}
\begin{equation}
  1\longrightarrow\mathbb Z/2\mathbb Z
  \xrightarrow{n}H\xrightarrow{f}G\longrightarrow e
  \tag{1.4.3.4}
\end{equation}
を記述する元とする。

\VerdierSemanticAnchor{ch1-1-4-3-cocycle}
$\alpha_s\in Z^2(G,\mathbb Z/2\mathbb Z)$ を、$f$ の正規化された切断 $s$
によって定まる、類 $\dot\alpha$ に属する正規化された非斉次余サイクルとする。
このとき

\VerdierSemanticAnchor{ch1-1-4-3-eq5}
\begin{equation}
  \begin{aligned}
    \alpha_s(\sigma_1,\sigma_2)
      &=s(\sigma_2)s(\sigma_1\sigma_2)^{-1}s(\sigma_1),
      & \sigma_1,\sigma_2&\in G,\\
    s(e)&=H\text{ の単位元},\\
    \alpha_s(\sigma_1,e)&=\alpha_s(e,\sigma_1)=1,\\
    \delta\alpha_s(\sigma_1,\sigma_2,\sigma_3)
      &=\alpha_s(\sigma_2,\sigma_3)
        \alpha_s(\sigma_1\sigma_2,\sigma_3)^{-1}\\
      &\quad{}\alpha_s(\sigma_1,\sigma_2\sigma_3)
        \alpha_s(\sigma_1,\sigma_2)^{-1}=1,
      & \sigma_1,\sigma_2,\sigma_3&\in G
  \end{aligned}
  \tag{1.4.3.5}
\end{equation}
である。

\VerdierSemanticAnchor{ch1-1-4-3-bijection}
切断 $s\colon G\to H$ は集合の同型
\[
  H\xrightarrow{\sim}\mathbb Z/2\mathbb Z\times G
\]
を定める。この同型によって $H$ の群法則を
$\mathbb Z/2\mathbb Z\times G$ に移すと、群法則

\VerdierSemanticAnchor{ch1-1-4-3-eq6}
\begin{equation}
  (\varepsilon_1,\sigma_1)(\varepsilon_2,\sigma_2)
  =\bigl(\varepsilon_1\varepsilon_2
    \alpha_s(\sigma_1,\sigma_2),\sigma_1\sigma_2\bigr)
  \tag{1.4.3.6}
\end{equation}
を得る。

\VerdierSemanticAnchor{ch1-1-4-3-identification}
以下 $H$ を (1.4.3.6) で記述された群と同一視する。任意の
$(\varepsilon_1,\sigma_1),(\varepsilon_2,\sigma_2)\in H$ に対して

\VerdierSemanticAnchor{ch1-1-4-3-eq7}
\begin{equation}
  \begin{aligned}
    T_f(\varepsilon_1,\sigma_1)&=T(\sigma_1),\\
    T'_f(\varepsilon_1,\sigma_1)&=T'(\sigma_1),\\
    c_f\bigl((\varepsilon_1,\sigma_1),(\varepsilon_2,\sigma_2)\bigr)
      &=c(\sigma_1,\sigma_2),\\
    c'_f\bigl((\varepsilon_1,\sigma_1),(\varepsilon_2,\sigma_2)\bigr)
      &=c'(\sigma_1,\sigma_2)
  \end{aligned}
  \tag{1.4.3.7}
\end{equation}
である。
